\documentclass[11pt]{article}

\usepackage[margin=1in]{geometry}
\usepackage[T1]{fontenc}
\usepackage[utf8]{inputenc}
\usepackage{lmodern}
\usepackage{booktabs}
\usepackage{enumitem}
\usepackage{parskip}
\usepackage[hidelinks]{hyperref}
\usepackage{xcolor}

\newcommand{\open}[1]{\textcolor{red!60!black}{\textbf{Open:} #1}}

\title{\textbf{bitsign MVP} \\ \large Engineering \& Product Design}
\author{Product and Engineering Draft}
\date{Draft v0.3 --- August 2026}

\begin{document}

\maketitle

\begin{abstract}
\noindent The bitsign MVP consists of a Rust backend deployed on Cloudflare, a
replaceable reference translation model, and an iOS application for near-real-time
sign-to-English inference. A separate, explicitly consented contribution workflow
collects structured feedback and training samples. The initial scope is narrow:
demonstrate an end-to-end communication path, test whether signing users find it
useful, and establish measurable quality, privacy, security, and accessibility
gates for further evaluation.
\end{abstract}

\tableofcontents

\section{Purpose and Context}

The MVP exists to prove three things, in order of importance:

\begin{enumerate}[nosep]
  \item \textbf{Demand} --- Deaf and signing users want this product and will use it.
  \item \textbf{The data loop} --- contributors will provide explicitly consented,
        rights-cleared samples and blind reviews, producing a continuously improving
        corpus for evaluation and model development.
  \item \textbf{The serving path} --- a working end-to-end demonstration (camera
        $\to$ inference $\to$ text/voice) meets defined performance, privacy, and
        accessibility gates.
\end{enumerate}

The translation model is replaceable from day one (\S\ref{sec:model}). Every model
version carries a capability label based on a frozen evaluation suite and ASL-fluent
human review. The initial label is \emph{technical prototype}.

\subsection{Release gates}

Two release gates can be prepared in parallel:

\begin{description}[nosep]
  \item[Technical evaluation.] The app and backend must pass end-to-end,
        performance, privacy, security, and accessibility tests with internal
        testers.
  \item[External evaluation.] The build must have a frozen capability label,
        ASL-fluent reviewer coverage, clear participant consent, and a documented
        stop condition before use with external participants.
\end{description}

External evaluation begins only after both gates pass.

\section{Product Surface}

Two modes sit behind one account and can be switched at any time:

\subsection{Communicate (user mode)}
Open app $\to$ camera captures signing $\to$ backend inference $\to$ meaning rendered
as text and natural voice for the hearing party. Reverse direction (speech $\to$ text
for the Deaf user to read) ships in the MVP build because it is cheap (on-device
speech recognition), low-risk, and completes a conversation loop; the community
evaluation covers only the sign-to-English direction, while the reply view remains
available. Signing avatars and multilingual routing are out of scope.

MVP capture model: short buffered clips (2--6\,s) submitted continuously, rendered as
a rolling transcript. True streaming inference is a post-MVP optimization; the UX
should feel live even if the implementation is clip-based.
Every transcript is displayed under the build's current capability label with a
correction affordance (thumbs up/down + optional corrected text). Corrections enter
evaluation data only after explicit consent.

\subsection{Contribute (labeler mode)}
Contribute has a separate consent flow stating purpose (model training and
evaluation), retention, access, withdrawal, and deletion rules. Communicate clips
are never training data by default. Contribute clips are training data by explicit
consent.

Two task types form a blind semantic round trip:

\begin{description}[nosep]
  \item[Task A --- meaning $\to$ sign.] Contributor receives a text prompt
        (``express this naturally in ASL''), records themselves signing, submits.
        The prompt is treated as a committed script: hashed at issuance, never shown
        to Task~B reviewers.
  \item[Task B --- sign $\to$ meaning.] Reviewer receives \emph{only} a video (never
        the prompt) and writes what it means. Each candidate clip is reviewed by
        $k$ independent reviewers ($k=3$ at MVP).
\end{description}

Consensus scoring on the $k$ blind reconstructions against the hidden script
determines the outcome (thresholds are placeholders, configurable server-side):
$\geq$85\% $\to$ accepted (clip enters the corpus as high-confidence labeled data);
70--84.9\% $\to$ variant/ambiguous (preserved for secondary review as candidate
regional-variation data); $<$70\% $\to$ flagged (investigate bad signing, bad
reconstruction, bad capture, or abuse; useful failures may be retained as edge or
adversarial data when consent permits).

Assignment of reviewers is server-side, randomized, and blind in both directions.
Contributors cannot review their own or their associates' clips (device/account
heuristics at MVP).

\section{Contribution Quality and Abuse Controls}

\begin{itemize}[nosep]
  \item \textbf{Acceptance is quality-gated.} Corpus admission depends on the blind
        review outcome, and reviewer accuracy is measured independently of whether
        a candidate clip passes.
  \item \textbf{Reviewer independence.} A reviewer cannot evaluate their own clip,
        repeated pairings are limited, and seeded known-answer tasks measure drift.
  \item \textbf{Anti-abuse floor at MVP:} one account per Apple ID, device
        attestation (App Attest), rate limits per account/device/IP, seeded
        known-answer tasks interleaved into review queues (anchor-clip pattern) to
        measure reviewer accuracy, and a manual review dashboard for suspicious
        activity.
\end{itemize}

\section{Contribution Data Design}

Every collected artifact must preserve enough provenance for reproducible
evaluation, model development, consent enforcement, and deletion:

\begin{itemize}[nosep]
  \item \textbf{Known-plaintext discipline.} Task prompts are scripts: stored with a
        commitment hash at issuance, encrypted at rest, revealed to scoring only
        after reconstructions are committed. The audit chain (script commit $\to$
        recording $\to$ video hash $\to$ blind assignment $\to$ reconstruction
        commit $\to$ match $\to$ manifest) is recorded in the database.
  \item \textbf{Reference sets.} Accepted reconstructions attach to the
        \emph{script} (not the clip) as paraphrase references, capped and
        tightness-checked. These references support deterministic evaluation while
        allowing valid language variation.
  \item \textbf{Consent and privacy.} Contribute's consent protocol
        covers training use explicitly. Communicate video is ephemeral (deleted
        after inference and excluded from training) unless the user explicitly opts
        a clip in. A de-identification pass runs before any corpus export. Deletion
        requests propagate to derived artifacts.
  \item \textbf{Rights provenance per clip:} script id and hash, contributor
        account, consent record version, capture timestamps, device class, reviewer
        set, consensus scores, and outcome. Skeletal pose extraction can run offline
        later; raw video is the artifact of record at MVP.
\end{itemize}

\section{Reference Model}
\label{sec:model}

The model slot is a replaceable interface. Every training artifact must have usage
rights appropriate for its intended deployment. Public sign-language corpora and
models carry varying restrictions, so each candidate requires a documented data and
model provenance review before it can leave internal evaluation. The long-term
training path uses the explicitly consented corpus produced by Contribute.

\begin{enumerate}[nosep]
  \item \textbf{Bring-up (internal only).} Evaluate a documented public reference
        model to stand up the serving path and evaluation harness. It carries the
        \emph{technical prototype} capability label and remains limited to internal
        evaluation until its provenance review is complete.
  \item \textbf{Isolated-sign path.} A small
        keypoint-based classifier (MediaPipe landmarks) for fingerspelling and a
        bounded vocabulary of isolated signs, trained on \emph{our own consented
        captures} --- the data requirement is small (26 handshapes + top-$N$ signs
        at 10 variants each is Contribute's first script set), and the model is fast
        enough to serve at the edge. It is the first candidate for the
        \emph{bounded vocabulary} capability label.
  \item \textbf{Continuous translation.} Fine-tune on the consented corpus as it
        grows. Capability labels advance only when frozen evaluation results support
        the change.
  \item \textbf{Model promotion.} A candidate replaces the active reference model
        only after it exceeds the current version on the held-out benchmark and
        passes the relevant release gates.
        The backend treats models as versioned artifacts with a manifest (weights
        hash, training-data rights class, benchmark scores, capability label);
        swapping models is an ops action, not a code change. The rights-class field
        prevents a restricted model from reaching an incompatible deployment
        surface.
\end{enumerate}

\open{benchmark candidate reference models on the frozen internal test set.}

\section{Backend Architecture}
\label{sec:backend}

Rust services use a shared Axum application pattern across two deployment planes.

\subsection{Control plane --- Cloudflare}

\begin{itemize}[nosep]
  \item \textbf{Worker front door} --- routing, auth verification, rate limiting,
        request metadata, and dispatch to application containers.
  \item \textbf{API service} --- Axum app on Cloudflare Containers via
        a platform-binding adapter: accounts, sessions, task issuance and
        assignment, consensus scoring, consent records, and admin/review endpoints.
  \item \textbf{Storage} --- R2 for video (upload via presigned URLs directly from
        the app; Communicate objects carry a TTL lifecycle rule implementing the
        ephemeral-processing guarantee); D1 for relational state (accounts, audit
        entries, scripts, tasks, reviews, consensus, model registry) reached through the
        command channel; KV for sessions/config/feature flags; Queues for
        asynchronous work (consensus evaluation, corpus export, deletion
        propagation, notification fan-out).
  \item \textbf{Audit log} is append-only, with idempotency keys on task issuance,
        submission, review, consent, and deletion mutations.
\end{itemize}

\subsection{Inference plane --- same code, GPU where needed}

Cloudflare Containers are CPU-class; the continuous VL model wants a GPU. Rather
than duplicate application logic, the inference service uses the same Axum interface
on a GPU-capable runtime:

\begin{itemize}[nosep]
  \item \textbf{Primary: Cloud Run GPU} --- the inference service is a second
        Axum app (clip in $\to$ transcript + confidence out), deployed to Cloud Run
        with an L4-class GPU and scale-to-zero during MVP. Planned evaluation
        sessions pre-warm the service to avoid cold-start delays.
  \item \textbf{Secondary: CPU path on Cloudflare Containers} for the
        isolated-sign/keypoint model (\S\ref{sec:model} item 2), which is light
        enough to run at the edge.
  \item The Worker front door routes inference requests to whichever plane serves
        the requested model version; the model registry records where each version
        lives and its rights class. \open{benchmark a quantized 2B VL model on CPU;
        if clip latency reaches p95 $\leq$ 12\,s end-to-end, defer the GPU plane.}
\end{itemize}

Technical release tests require strict input and media validation with stable errors
(no placeholder output); bounded media size, duration, decoded frames,
generated-text length, wall time, and concurrency; deletion tests across every
storage layer; and telemetry that excludes video, frames, returned text, and
identity.

\subsection{API sketch (MVP)}

\texttt{POST /v1/auth/apple} $\cdot$
\texttt{POST /v1/infer} (clip $\to$ transcript + capability label)
$\cdot$ \texttt{GET /v1/tasks/next} $\cdot$
\texttt{POST /v1/tasks/\{id\}/submit} (Task A upload) $\cdot$
\texttt{POST /v1/reviews/\{id\}/submit} (Task B reconstruction) $\cdot$
\texttt{GET /v1/models/current} $\cdot$
\texttt{DELETE /v1/me} (account + data deletion) $\cdot$
\texttt{GET /healthz}.

\section{iOS Application}

SwiftUI, iOS 17+, distributed through TestFlight for controlled evaluation.

\begin{itemize}[nosep]
  \item \textbf{Screens (MVP):} onboarding + capability-label introduction,
        Sign in with Apple, Communicate (camera with framing guide, clip review
        before upload, rolling transcript, voice output, speech$\to$text reply
        view), Contribute (study consent, task list, capture flow, review flow),
        and Settings (consent status, versions, deletion).
  \item \textbf{Capture:} AVFoundation, front camera, 720p clips, on-device
        segmentation into 2--6\,s chunks with a hard 6\,s stop; review-before-upload;
        background upload to R2 via presigned URLs; distinct upload vs.\ inference
        progress; graceful degradation on poor connectivity (Contribute works fully
        offline-queue, Communicate does not).
  \item \textbf{Voice output:} AVSpeechSynthesizer at MVP (natural TTS via a
        provider is a post-MVP swap). Speech input: SFSpeechRecognizer on-device.
  \item \textbf{Accessibility is a release gate:} VoiceOver labels and reading
        order, Dynamic Type, contrast, and a visual equivalent for every
        audio-communicated status.
  \item \textbf{Trust \& safety surface:} capability label on every transcript,
        correction affordance (feeds consented feedback), report/flag on any task
        content, ``aid, not interpreter replacement'' framing in onboarding.
\end{itemize}

\section{Milestones}

The order is chosen so each milestone ends with a testable increment. Community
definition, reviewer recruitment, and frozen-evaluation preparation can proceed in
parallel.

\begin{center}
\begin{tabular}{@{}llp{7.6cm}@{}}
\toprule
\# & Target & Definition of done \\
\midrule
M1 & wk 1--2 & Backend skeleton live on Cloudflare: auth, consent records, audit log,
R2 upload with TTL lifecycle, and \texttt{/healthz} at a documented endpoint.
Bring-up model evaluated (\S\ref{sec:model}). \\
M2 & wk 2--4 & Inference plane serving the bring-up model behind the registry;
\texttt{POST /v1/infer} returns a transcript for a real clip with stable errors on
invalid media. \\
M3 & wk 3--5 & iOS Communicate path against live backend; capability label shown;
TestFlight build installable by an internal tester. \\
M4 & wk 5--6 & Contribute pipeline end-to-end under its consent protocol (Task A
$\to$ blind Task B $\times$3 $\to$ consensus); first consented script set
(fingerspelling + top-$N$ signs) live. \\
M5 & wk 6 & Recorded end-to-end demonstration: a signed request from the TestFlight
build returns an inference result from the live backend. \\
M6 & post-MVP & Frozen ASL-fluent evaluation of the isolated-sign model for the
\emph{bounded vocabulary} label, followed by a closed external evaluation after all
release gates pass. \\
\bottomrule
\end{tabular}
\end{center}

Weeks are effort estimates rather than calendar commitments.

\section{Non-Goals (MVP)}

Public production release; production hardening and high availability; general
feature development beyond the demonstration path; avatar rendering; multilingual
routing; Android; web app; natural-voice cloning; enterprise features; and external
evaluation before the release gates pass.

\section{Open Questions}

\begin{enumerate}[nosep]
  \item Candidate reference-model evaluation on the frozen internal clip set
        (\S\ref{sec:model}).
  \item CPU-only inference viability --- decides whether the Cloud Run GPU plane
        ships in MVP (\S\ref{sec:backend}).
  \item First consented script-set contents: fingerspelling plus which top-$N$ signs.
  \item ASL-fluent reviewer coverage and community-adviser recruitment for external
        evaluation.
  \item TestFlight tester list for technical evaluation.
  \item Final bitsign wordmark and in-app capability-label presentation.
  \item Corpus export layout: R2 bucket structure and manifest format.
  \item Verification of external evidence supporting capability and safety claims
        before publication.
\end{enumerate}

\end{document}
